% Independent mathematical derivation. Prefix HTP. Requires amsmath,amssymb,hyperref.
\section{The exact product transported by the original heat group}
\label{sec:heat-transported-product}

\paragraph{Sources and scope.}
The analytic domain and heat operator are the original HSL15--18
in the complete included \nolinkurl{FULL_LATTICE_HEAT_OPERATOR.tex}.
Lines 366--469 were read directly at source SHA--256
\begin{center}\small
\nolinkurl{42c1de0bb8991f460f51d8f84c31d271b01fa5f5b9e6a7f2e803a0ea6deb6841}.
\end{center}
The norm estimates needed below are proved here as well.
The actual theta source is Brad Rodgers and Terence Tao,
\href{https://arxiv.org/abs/1801.05914v5}{\emph{The De Bruijn--Newman
constant is non-negative}}, author file \nolinkurl{fmp-template.tex},
SHA--256
\begin{center}\small
\nolinkurl{25b015c0fb151f1613247d82cf53a6561a320d4dab2bb4f809482e7c477a5817}.
\end{center}
The prior actual reading is lines 1--155, including
\texttt{phidef} and \texttt{htdef}; no whole-paper reading is
claimed. The positive-kernel proof and all heat integrals are
THS1--29, proved and read completely in the included
\nolinkurl{THETA_HEAT_FULL_SUPPORT.tex}. Its current source
SHA--256, including the later tail bounds THS30--31, is
\begin{center}\small
\nolinkurl{69507dcc5e192fad2023fff6e7fac1f435dbb2dd1809e6c68d1b5b1e712a9957}.
\end{center}
No threshold theorem or claim about the location of the theta
family's zeros is a premise. This is the amplitude-ring
calculation; the independent coefficient supports and their
full lattice map remain the subject of the included supported
heat construction, rather than being identified with one label.

\subsection{The complete analytic domain and its exact norm bounds}

For $h>0$, retain precisely the coefficient norms
\begin{equation}\tag{HTP1}
 \|f\|_h=\sum_{n\ge0}|a_n|\sqrt{n!}\,h^n,
 \qquad f(z)=\sum_{n\ge0}a_nz^n,\qquad
 \mathcal E=\{f:\|f\|_h<\infty\text{ for every }h>0\}.
\end{equation}
The positive integer $h$ give a countable defining family.
A Cauchy sequence in every such norm has coordinate limits;
completeness of weighted $\ell^1$ gives the same limit in each
norm. Thus $\mathcal E$ is a complete Fr\'echet space.
Every polynomial truncation converges in all the norms, since
each omitted weighted tail is summable. The coefficient domain
equals the entire-function domain
\begin{equation}\tag{HTP2}
\begin{gathered}
 p_\alpha(f)=\sup_{z\in\mathbb C}|f(z)|e^{-\alpha|z|^2},\qquad
 \mathcal E=\{f\text{ entire}:p_\alpha(f)<\infty\ (\alpha>0)\},\\
 p_{1/(2h^2)}(f)\le\|f\|_h,\qquad
 \|f\|_h\le\frac{p_\alpha(f)}{1-h\sqrt{2\mathrm e\alpha}}
       \quad(h\sqrt{2\mathrm e\alpha}<1).
\end{gathered}
\end{equation}
Here $\mathrm e$ denotes the ordinary exponential constant.
For the first inequality, the square of
$|z|^n/(h^n\sqrt{n!})$ is a term of
$e^{|z|^2/h^2}$. Multiply the resulting bound by the weighted
coefficients and sum. This proves entire convergence too.
Conversely, Cauchy's coefficient inequality on radius
$r=\sqrt{n/(2\alpha)}$, for $n\ge1$, and $n!\le n^n$ give
$|a_n|\sqrt{n!}h^n\le
p_\alpha(f)(h\sqrt{2\mathrm e\alpha})^n$.
The degree-zero term is at most $p_\alpha(f)$; the geometric
sum proves the second inequality and the stated equivalence.

Ordinary multiplication is a continuous bilinear operation:
\begin{equation}\tag{HTP3}
                \|fg\|_h\le\|f\|_{\sqrt2h}\|g\|_{\sqrt2h}.
\end{equation}
Indeed, for coefficient indices $j,k\ge0$,
$\sqrt{(j+k)!}\le\sqrt{j!k!}\,2^{(j+k)/2}$, since
$\binom{j+k}{j}\le2^{j+k}$. Taking absolute values before
the Cauchy-product sum yields HTP3 and its absolute convergence.
For every integer $n\ge0$ and every $b>0$ one also has
\begin{equation}\tag{HTP4}
          \|D^nf\|_h\le\sqrt{n!}\,b^{-n}
                                \|f\|_{\sqrt{h^2+b^2}},
          \qquad D=\frac{d}{dz}.
\end{equation}
For $k\ge n$, the contribution of $a_k$ to the left side is
$|a_k|\sqrt{k!}\sqrt{n!}\sqrt{\binom kn}h^{k-n}$.
The single nonnegative binomial term gives
$\binom kn h^{2(k-n)}b^{2n}\le(h^2+b^2)^k$.
Its square root proves the bound after summation.
Thus every derivative is a continuous operator on this exact
domain, with its original factorials.

The heat group is
\begin{equation}\tag{HTP5}
\begin{gathered}
 T_tf=\sum_{m\ge0}\frac{(-t)^m}{m!}D^{2m}f,\qquad
 (T_tf)_i=\sum_{m\ge0}\frac{(-t)^m(i+2m)!}{m!i!}a_{i+2m},\\
 \sum_{m\ge0}\frac{|t|^m}{m!}\|D^{2m}f\|_h
                       \le\|f\|_{h+\sqrt{2|t|}},\qquad t\in\mathbb C.
\end{gathered}
\end{equation}
To prove the last inequality, put $k=i+2m$ in the nonnegative
sum of coefficient norms. Its inner factor is
\[
 \sum_{m=0}^{\lfloor k/2\rfloor}
 \frac{\sqrt{k!}}{m!\sqrt{(k-2m)!}}|t|^m h^{k-2m}.
\]
The numerical coefficient equals
$\sqrt{\binom k{2m}}\sqrt{\binom{2m}m}$, at most
$\binom k{2m}2^m$. The sum is bounded by the even part of
$(h+\sqrt{2|t|})^k$, and therefore by that whole expression.
Summing against $|a_k|\sqrt{k!}$ proves HTP5.
Every heat sum is thus absolutely convergent in every defining
norm, locally uniformly in $t$ by enlarging $|t|$.
Moreover
\begin{equation}\tag{HTP6}
 T_tT_s=T_{t+s},\qquad T_t^{-1}=T_{-t},\qquad T_0=I,
 \qquad \partial_tT_tf=-D^2T_tf.
\end{equation}
The absolute double sum is at most
$\|f\|_{h+\sqrt{2|t|}+\sqrt{2|s|}}$ by two applications of
HTP5. Regrouping by total derivative order is therefore
permitted. Its finite binomial coefficients give $t+s$.
The identity at zero and inverse follow. For differentiation,
bound the series on a larger time disc by HTP5 and use the
convergent power series in the complete normed coordinate
spaces. This gives the displayed derivative in every norm.

\subsection{A convergent product, with the exact derivative formula}

For $a,b>0$ satisfying $ab>2|t|$, HTP3--4 prove
\begin{equation}\tag{HTP7}
 \sum_{n\ge0}\frac{(2|t|)^n}{n!}\|f^{(n)}g^{(n)}\|_h
 \le\frac{\|f\|_{\sqrt{2h^2+a^2}}
             \|g\|_{\sqrt{2h^2+b^2}}}{1-2|t|/(ab)}.
\end{equation}
Apply HTP4 with parameters $a,b$ at the norm $\sqrt2h$ to
the two derivative factors in HTP3. The product contributes
$n!(ab)^{-n}$, which cancels the $n!$ in the denominator.
The remaining geometric series has the stated sum. Every
factor and radius in HTP7 is therefore explicit. The bound
is uniform on a compact time set after choosing $ab$ greater
than twice its largest $|t|$. It proves absolute convergence
in $\mathcal E$ of the derivative-product series.

Define the transported product by the actual continuous heat
automorphism of the vector space:
\begin{equation}\tag{HTP8}
\begin{gathered}
 f\star_tg=T_t\bigl((T_{-t}f)(T_{-t}g)\bigr),\qquad
 \mathcal E_t=(\mathcal E,\star_t),\\
 \|f\star_tg\|_h\le\|f\|_{R_t(h)}\|g\|_{R_t(h)},\qquad
 R_t(h)=\sqrt2\bigl(h+\sqrt{2|t|}\bigr)+\sqrt{2|t|}.
\end{gathered}
\end{equation}
First apply HTP5 to the outer heat operator, then HTP3 at
$h+\sqrt{2|t|}$, and then HTP5 to each inner operator.
This proves the norm bound and bilinear continuity directly.

For polynomials the product has a finite operator calculation.
On a polynomial $f(z_1)g(z_2)$ put $D_i=\partial_{z_i}$ and
let $\mu$ set $z_1=z_2=z$. Leibniz's rule gives
$D\mu=\mu(D_1+D_2)$. Every differential exponential below
terminates on this polynomial; the differential operators
commute. Consequently
\begin{equation}\tag{HTP9}
 \begin{split}
 T_t\bigl((T_{-t}f)(T_{-t}g)\bigr)
 &=\mu e^{-t(D_1+D_2)^2}e^{t(D_1^2+D_2^2)}(f\otimes g)\\
 &=\mu e^{-2tD_1D_2}(f\otimes g)
   =\sum_{n\ge0}\frac{(-2t)^n}{n!}f^{(n)}g^{(n)}.
 \end{split}
\end{equation}
Combining commuting finite exponentials follows by the finite
binomial formula at each order, not by cancelling divergent
series. Now approximate $f,g\in\mathcal E$ by their polynomial
truncations in every norm. HTP8 makes the left side continuous
in both arguments. HTP7 makes the sum on the right a continuous
bilinear map: apply that bound to the difference in each
argument, and note that a convergent sequence is bounded in
each input norm. Thus the polynomial identity extends to
\begin{equation}\tag{HTP10}
 \boxed{\displaystyle
 f\star_tg=\sum_{n=0}^{\infty}\frac{(-2t)^n}{n!}
                                      f^{(n)}g^{(n)}.}
\end{equation}
This proof establishes the identity on the exact space, with
absolute convergence, without a formal infinite cancellation.

Even the expansion of all three heat operators and every
Leibniz term is absolutely bounded. For $x=|t|$ and
$R=\sqrt2(h+\sqrt{2x})+\sqrt{2x}$ one has
\begin{equation}\tag{HTP11}
 \sum_{a,b,c\ge0}\sum_{v=0}^{2a}
 \frac{x^{a+b+c}}{a!b!c!}\binom{2a}{v}
 \bigl\|f^{(2b+v)}g^{(2c+2a-v)}\bigr\|_h
                         \le\|f\|_R\|g\|_R.
\end{equation}
Let $f^\#(z)=\sum|a_n|z^n$ and define $g^\#$ similarly.
They belong to $\mathcal E$ with the same norms. Derivatives
and products of these series have nonnegative coefficients;
they dominate coefficientwise the absolute coefficients of
the corresponding derivatives and products of $f,g$.
With $S_x=T_{-x}$, the nonnegative coefficient expansion of
$S_x((S_xf^\#)(S_xg^\#))$ has exactly the weights in HTP11,
by Leibniz's rule. For any finite set of terms its coefficient
sum is bounded by this entire nonnegative series. The norm of
the latter is at most $\|f\|_R\|g\|_R$ by HTP8's proof
using three positive heat sums. Monotone convergence in the
nonnegative coefficient sums proves HTP11. Thus expansions
or regroupings of this product, when desired, have a separate
finite absolute majorant.

The product is commutative and associative and has unit $1$:
\begin{equation}\tag{HTP12}
\begin{gathered}
 T_{-t}(f\star_tg)=(T_{-t}f)(T_{-t}g),\qquad
 (f\star_tg)\star_tk=f\star_t(g\star_tk),\qquad
 f\star_t1=f,\\
 T_t:(\mathcal E,\cdot)\overset{\sim}{\longrightarrow}
       (\mathcal E,\star_t)
 \quad\text{is a continuous unital algebra isomorphism.}
\end{gathered}
\end{equation}
The first identity is HTP6 applied to the definition. Applying
the injective map $T_{-t}$ to either parenthesization gives the
same ordinary product of three functions; this proves
associativity with no triple derivative-series rearrangement.
Commutativity follows the same way. Since $T_t1=1$, the unit
claim follows too. Bilinearity and both directions of continuity
were already proved. The rings have the same vector-space
zero and unit, while their products are the specified ones.

\subsection{Characters, principal ideals and the exact finite jets}

For $t\ne0$, ordinary evaluation at any point fails to be a
character of this product. With $w$ denoting the coordinate
function,
\begin{equation}\tag{HTP13}
             w\star_tw=w^2-2t,
 \qquad \operatorname{ev}_z(w\star_tw)=z^2-2t
                          \ne z^2=\operatorname{ev}_z(w)^2.
\end{equation}
Only $n=0,1$ contribute in HTP10; the factor $2t$ and its
minus sign are therefore fixed.

We first establish the division needed to characterize all
complex-valued characters, without assuming their continuity.
For every $f\in\mathcal E$ and $z\in\mathbb C$, the removable
quotient belongs to the same space:
\begin{equation}\tag{HTP14}
 Q_zf(w)=\begin{cases}
       (f(w)-f(z))/(w-z),&w\ne z,\\ f'(z),&w=z,
       \end{cases}
 \qquad Q_zf\in\mathcal E,
 \qquad f-f(z)=(w-z)Q_zf.
\end{equation}
For $|w-z|\ge1$ its $p_\alpha$-weighted absolute value is at
most $p_\alpha(f)+|f(z)|\le
p_\alpha(f)(1+e^{\alpha|z|^2})$.
For $|w-z|\le1$, use
$Q_zf(w)=\int_0^1f'(z+s(w-z))\,ds$. Cauchy's derivative
bound with radius $1$ at each point of this segment is at
most $p_\alpha(f)e^{\alpha(|z|+2)^2}$.
Multiplication by $e^{-\alpha|w|^2}\le1$ preserves that
bound. This proves finite $p_\alpha(Q_zf)$ for every
$\alpha>0$, hence HTP14 by HTP2. Repeated division proves
the corresponding finite-order Taylor remainder statement.

Every unital $\mathbb C$-algebra character on the ordinary
ring $\mathcal E$ is evaluation. Indeed, for a character
$\chi$ put $z=\chi(w)$. Applying $\chi$ to HTP14 gives
$\chi(f)=f(z)$, because the factor $\chi(w-z)$ is zero.
Conversely evaluation is visibly a unital character and is
continuous by HTP2. HTP12 now gives the complete transported
character formula
\begin{equation}\tag{HTP15}
 \chi_{t,z}(f)=(T_{-t}f)(z),\qquad z\in\mathbb C,
 \qquad
 \operatorname{Hom}_{\mathbb C\text{-alg},1}(\mathcal E_t,\mathbb C)
                  =\{\chi_{t,z}:z\in\mathbb C\}.
\end{equation}
For a character on $\mathcal E_t$, compose with $T_t$ and
apply the ordinary classification just proved. Distinct $z$
give distinct characters because $T_{-t}w=w$. The formulas
and HTP5/HTP2 show that every one is continuous. No statement
about maximal ideals with other residue fields is needed.

For any $F\in\mathcal E$, exact principal-ideal transport is
\begin{equation}\tag{HTP16}
\begin{gathered}
 T_t(F\mathcal E)=(T_tF)\star_t\mathcal E,\\
 \mathcal E/(F\mathcal E)\overset{\sim}{\longrightarrow}
 \mathcal E_t/((T_tF)\star_t\mathcal E),\qquad
 [g]\longmapsto[T_tg].
\end{gathered}
\end{equation}
If $g=Fq$ then $T_tg=(T_tF)\star_t(T_tq)$ by HTP12.
Conversely every element on the right has this form because
$T_t$ is onto. Its inverse $T_{-t}$ proves the equality of
ideals and the well-defined, injective and surjective quotient
map, preserving both operations and both identities.
This is an algebraic ideal equality; no closure of the ideal
is substituted for it.

For a point $z$ and integer $m\ge1$, write
$\mathfrak m_z=(w-z)\mathcal E$ and
$\mathfrak m_{t,z}=\ker\chi_{t,z}$. Then
\begin{equation}\tag{HTP17}
\begin{gathered}
 \mathfrak m_{t,z}=T_t\mathfrak m_z=(w-z)\star_t\mathcal E,
 \qquad (w-z)\star_tf=(w-z)f-2tf',\\
 T_t((w-z)^m)=
 \sum_{b=0}^{\lfloor m/2\rfloor}
       \frac{(-t)^bm!}{b!(m-2b)!}(w-z)^{m-2b}
       =(w-z)^{\star_t m},\\
 \mathfrak m_{t,z}^{\star_t m}
                 =T_t((w-z)^m\mathcal E).
\end{gathered}
\end{equation}
The first identity follows from HTP14--16 and $T_t(w-z)=w-z$.
The derivative formula has only two terms in HTP10.
The polynomial formula is the finite heat series HTP5.
The algebra isomorphism transports powers and products of
ideals, proving the last two assertions with no formal
interpretation of the displayed polynomial required.

The exact jet map, including its inverse on representatives, is
\begin{equation}\tag{HTP18}
\begin{gathered}
 J_{t,z,m}(f)=\sum_{k=0}^{m-1}
       \frac{(D^kT_{-t}f)(z)}{k!}\epsilon^k
           \in\mathbb C[\epsilon]/(\epsilon^m),\\
 \ker J_{t,z,m}=\mathfrak m_{t,z}^{\star_t m},\qquad
 \mathcal E_t/\mathfrak m_{t,z}^{\star_t m}
                \overset{\sim}{\longrightarrow}
                   \mathbb C[\epsilon]/(\epsilon^m),\\
 \sum_{k<m}c_k\epsilon^k\longmapsto
       \left[T_t\left(\sum_{k<m}c_k(w-z)^k\right)\right].
\end{gathered}
\end{equation}
The ordinary Taylor map is multiplicative modulo
$\epsilon^m$ by the finite Leibniz formula and preserves
zero and unit. Its kernel consists precisely of functions
divisible by $(w-z)^m$ in $\mathcal E$: apply HTP14
successively to the zero Taylor remainders. Its image contains
every polynomial of degree less than $m$ in $\epsilon$.
Compose with the isomorphism $T_{-t}$ to prove all lines of
HTP18. In particular the transported multiplicity of $T_tF$
at $\chi_{t,z}$ is the original order of $F$ at $z$.
This is an equality of the specified finite jets; it is not
an equality with the ordinary Taylor jets of $T_tF$.

\subsection{The actual theta family and its complete double integral}

Keep the original source kernel, with no change in variables:
\begin{equation}\tag{HTP19}
\begin{gathered}
 \Phi(u)=\sum_{n\ge1}
       (2\pi^2n^4e^{9u}-3\pi n^2e^{5u})e^{-\pi n^2e^{4u}},
 \quad u\ge0,\\
 H_a(z)=\int_0^\infty e^{au^2}\Phi(u)\cos(zu)\,du,\qquad a\in\mathbb C,\\
 0<\pi(2\pi-3)e^{9u-\pi e^{4u}}
       \le\Phi(u)\le C_\Phi e^{9u-\pi e^{4u}},\\
 C_\Phi=2\pi^2\frac{1+11c+11c^2+c^3}{(1-c)^5},\qquad c=e^{-3\pi}.
\end{gathered}
\end{equation}
For completeness, every original summand is positive because
its prefactor is $\pi n^2e^{5u}(2\pi n^2e^{4u}-3)$.
The $n=1$ term gives the lower bound using
$2\pi e^{4u}-3\ge(2\pi-3)e^{4u}$.
For the upper bound discard the negative term in each
prefactor and use $n^2-1\ge3(n-1)$ and $e^{4u}\ge1$.
The remaining sum is
$\sum_{n\ge1}n^4c^{n-1}=(1+11c+11c^2+c^3)/(1-c)^5$;
four differentiations of the geometric series give this
identity. This proves the entire stated bound directly.

The full moment control, retaining all constants, is
\begin{equation}\tag{HTP20}
\begin{gathered}
 I_p(A,B):=\int_0^\infty u^pe^{Au^2+Bu}\Phi(u)\,du
 \le\mathcal M_p(A,B),\\
 \mathcal M_p(A,B):=
 \frac{C_\Phi}{2\pi}
 \exp\left(\frac{3A^2}{32\pi}
       +\frac{(B+9+p)^2}{8\pi}-\frac\pi2\right),
 \quad A,B\ge0,\ p\ge0,\\
 p_\alpha(H_a)\le I_0(\max(\operatorname{Re}a+1/(4\alpha),0),0),
 \qquad H_a\in\mathcal E,\qquad T_rH_a=H_{a+r}.
\end{gathered}
\end{equation}
The moment proof uses $u^p\le e^{pu}$ and splits the
negative kernel exponent into quarters and a half. Namely,
\[
 Au^2-\frac\pi4e^{4u}\le\frac{3A^2}{32\pi},\qquad
 (B+9+p)u-\frac\pi4e^{4u}\le\frac{(B+9+p)^2}{8\pi}.
\]
The first follows from $e^{4u}\ge(32/3)u^4$ and completing
the square in $u^2$; the second follows from $e^{4u}\ge8u^2$
and completing the square in $u$.
The remaining integral is at most $e^{-\pi/2}/(2\pi)$,
because $e^{4u}\ge1+4u$. This proves the first line.
For the second, $|\cos(zu)|\le e^{|z|u}$ and
$|z|u\le\alpha|z|^2+u^2/(4\alpha)$ prove the growth bound,
so HTP2 gives $H_a\in\mathcal E$.
Differentiating the cosine under the integral is justified
at every order by $I_p$ on compact parameter sets. The sum
of the absolute heat-integrands is bounded by
$e^{(\operatorname{Re}a+|r|)u^2+|\operatorname{Im}z|u}\Phi(u)$,
with its quadratic coefficient enlarged to zero if negative.
It is integrable by the first line. Expanding $e^{ru^2}$
therefore proves $T_rH_a=H_{a+r}$ with the original heat
sign; equality of entire functions identifies the same
elements of $\mathcal E$. No assertion concerning their
zero locations occurs in this argument.

For arbitrary complex $a,b,t,z$, the actual transported
product is the following exact double integral:
\begin{equation}\tag{HTP21}
\boxed{\begin{aligned}
 (H_a\star_tH_b)(z)
  =\frac12\int_0^\infty\!\int_0^\infty
       &\Phi(u)\Phi(v)e^{au^2+bv^2}\\[-2pt]
       &\cdot\bigl[e^{2tuv}\cos(z(u+v))
                    +e^{-2tuv}\cos(z(u-v))\bigr]\,du\,dv.
\end{aligned}}
\end{equation}
All integrals, derivative sums and their exchanges converge
in every norm of $\mathcal E$, as we now verify explicitly.
For $x\ge0$,
$\sum_{k\ge0}x^k/\sqrt{k!}\le\sqrt2e^{x^2}$:
write each summand as
$2^{-k/2}(\sqrt2x)^k/\sqrt{k!}$ and apply Cauchy--Schwarz.
Thus both the sine and cosine of $uz$, for real $u$, have
$h$-norm at most $\sqrt2e^{h^2u^2}$.
Since a derivative of order $n$ changes only sine to cosine
and supplies a factor $u^n$, HTP3 gives the majorant
\begin{equation}\tag{HTP22}
\begin{gathered}
 A_h=\max(\operatorname{Re}a+2h^2+|t|,0),\quad
 B_h=\max(\operatorname{Re}b+2h^2+|t|,0),\\
 \int_0^\infty\!\int_0^\infty
 \Phi(u)\Phi(v)e^{(\operatorname{Re}a)u^2+(\operatorname{Re}b)v^2}
 \sum_{n\ge0}\frac{(2|t|)^n}{n!}
 \bigl\|D^n\cos(uz)D^n\cos(vz)\bigr\|_h\,du\,dv\\
 \hspace{20mm}\le2I_0(A_h,0)I_0(B_h,0)
       \le2\mathcal M_0(A_h,0)\mathcal M_0(B_h,0)<\infty.
\end{gathered}
\end{equation}
Indeed the norm of the derivative product is at most
$2(uv)^ne^{2h^2(u^2+v^2)}$. Its weighted sum is
$2e^{2|t|uv+2h^2(u^2+v^2)}$, and
$2|t|uv\le|t|(u^2+v^2)$ separates the variables.
For the integrals of the individual derivatives use the
same bounds with the finite extra moments in HTP20.
Nonnegative summation of the absolute coefficients now
justifies integrating every derivative, multiplying the
two integrals and interchanging the derivative series with
the double integral in each $h$-norm. The resulting
coefficient integrals give a member of every weighted
$\ell^1$ space, hence the same member of $\mathcal E$.

For the value of that sum, the even derivative orders give
$\cos(zu)\cos(zv)\cosh(2tuv)$, and the odd orders give
$-\sin(zu)\sin(zv)\sinh(2tuv)$. The elementary cosine sum
and difference identities transform their sum into the
bracket in HTP21 divided by $2$. This proves that formula
with its signs. Alternatively its two cosine frequencies
are eigenvectors of $T_t$ with respective multipliers
$e^{t(u+v)^2}$ and $e^{t(u-v)^2}$; the already proved
absolute bounds justify the same calculation. On compact
sets of the complex parameters, HTP22 supplies one common
bound by enlarging their real parts and moduli. Additional
parameter derivatives contribute finite powers of $u,v$
and are bounded by the corresponding $I_p$ products.

The requested actual-time specialization and its transport
relation are therefore
\begin{equation}\tag{HTP23}
\begin{gathered}
 H_t\star_tH_s=T_t(H_0H_{s-t}),\qquad
 H_t\star_tH_t=T_t(H_0^2),\\
 (H_t\star_tH_s)(z)=\frac12\int_0^\infty\!\int_0^\infty
 \Phi(u)\Phi(v)e^{tu^2+sv^2}
   \bigl[e^{2tuv}\cos(z(u+v))+e^{-2tuv}\cos(z(u-v))\bigr]\,du\,dv.
\end{gathered}
\end{equation}
Use HTP20 in the definition HTP8, or set $a=t,b=s$ in
HTP21. For real $a,b,t$, the value of HTP21 at $z=0$ is
strictly positive: its bracket is $2\cosh(2tuv)>0$ and
the kernel factors are positive. That statement is about
this value of the actual integral; it supplies no assertion
that the entire function has no zeros at other points.

Finally the character and divisor comparison for the actual
family is exact:
\begin{equation}\tag{HTP24}
\begin{gathered}
 \chi_{t,z}(H_a)=H_{a-t}(z),\qquad
 \chi_{t,z}(H_t)=H_0(z),\qquad H_t(z)=\chi_{t,z}(H_{2t}),\\
 T_t(H_0\mathcal E)=H_t\star_t\mathcal E,\qquad
 \mathcal E/(H_0\mathcal E)
     \simeq\mathcal E_t/(H_t\star_t\mathcal E),\\
 J_{t,z,m}(H_t)=\sum_{k=0}^{m-1}\frac{H_0^{(k)}(z)}{k!}\epsilon^k,
 \qquad
 J_{t,z,m}(H_{2t})=\sum_{k=0}^{m-1}\frac{H_t^{(k)}(z)}{k!}\epsilon^k.
\end{gathered}
\end{equation}
The first line is HTP15 and HTP20. Apply HTP16 with
$F=H_0$ for the second, and HTP18 for the third.
The first transported principal divisor thus has the
original $H_0$ character zeros and original multiplicities.
The ordinary analytic zeros and multiplicities of $H_t$
are instead those of $H_{2t}$ tested by these transported
characters and jets. This supplies the exact comparison
between both objects; ordinary evaluation cannot be
substituted for a transported character when $t\ne0$ by
HTP13. None of these isomorphisms decides the zero-location
question for $H_0$ or asserts a conclusion about RH.

\begin{figure}[ht]
\centering
$\begin{array}{ccc}
 (\mathcal E,\cdot)&\xrightarrow{\quad T_t\quad}&(\mathcal E,\star_t)\\[3pt]
 {\scriptstyle\operatorname{ev}_z}\downarrow&&
                       \downarrow{\scriptstyle\chi_{t,z}=\operatorname{ev}_zT_{-t}}\\[3pt]
 \mathbb C&\xrightarrow{\quad\operatorname{id}\quad}&\mathbb C
 \end{array}$
\caption{The exact algebra and character comparison, HTP12--18.
The top map sends $H_0$ to $H_t$, its principal ideal to the
corresponding star ideal, and each finite jet to HTP18. Thus
the right character evaluates $H_t$ as $H_0(z)$.
The original analytic value $H_t(z)$ instead equals
$\chi_{t,z}(H_{2t})$, HTP24. The diagram uses the actual
Rodgers--Tao kernel HTP19 and the full proved heat action;
it does not identify ordinary evaluation with the right
character when $t\ne0$.}
\end{figure}
